% ============================================================
% SOFTWARE DESIGN
% ============================================================

\section{Software Design}

\subsection{Software Modules and Responsibilities}

Identify the main software modules and describe each module's responsibility,
inputs, outputs, and interfaces. Explain how the modules are organized across
controllers, applications, or services, as applicable.

The example below illustrates how input validation, processing, decision logic,
and fault handling can be separated into functions. Replace it with
representative project-specific code and select the appropriate language;
these function names and this structure are illustrative, not mandatory.

\begin{lstlisting}[language=Python, caption=Illustrative Software Module Structure]
# Template example only: replace with project-specific logic.
# Helper functions below must be defined for the actual project.

def process_record(record):
    if not validate_record(record):
        return handle_fault("INVALID_INPUT")

    processed_data = process_input(record)
    decision = evaluate_conditions(processed_data)
    return build_status(decision)
\end{lstlisting}

\subsection{Data Flow and Communication}

Describe how data moves between software modules and connected devices.
Define the relevant message formats, commands, sensor records, status fields,
and communication protocols. Record timing requirements and data validation
where applicable. Refer to the operating algorithm in the methodology section
rather than repeating it here.

\subsection{Fault Handling and Recovery}

Describe how software detects, logs, and communicates faults. Define the
responsible modules, fault states, recovery actions, and conditions for
resuming operation. Relate these mechanisms to the system failure recovery
protocol without repeating the system-level description.
